% Options for packages loaded elsewhere
\PassOptionsToPackage{unicode}{hyperref}
\PassOptionsToPackage{hyphens}{url}
\documentclass[
  french,
]{article}
\usepackage{xcolor}
\usepackage{amsmath,amssymb}
\setcounter{secnumdepth}{-\maxdimen} % remove section numbering
\usepackage{iftex}
\ifPDFTeX
  \usepackage[T1]{fontenc}
  \usepackage[utf8]{inputenc}
  \usepackage{textcomp} % provide euro and other symbols
\else % if luatex or xetex
  \usepackage{unicode-math} % this also loads fontspec
  \defaultfontfeatures{Scale=MatchLowercase}
  \defaultfontfeatures[\rmfamily]{Ligatures=TeX,Scale=1}
\fi
\usepackage{lmodern}
\ifPDFTeX\else
  % xetex/luatex font selection
\fi
% Use upquote if available, for straight quotes in verbatim environments
\IfFileExists{upquote.sty}{\usepackage{upquote}}{}
\IfFileExists{microtype.sty}{% use microtype if available
  \usepackage[]{microtype}
  \UseMicrotypeSet[protrusion]{basicmath} % disable protrusion for tt fonts
}{}
\makeatletter
\@ifundefined{KOMAClassName}{% if non-KOMA class
  \IfFileExists{parskip.sty}{%
    \usepackage{parskip}
  }{% else
    \setlength{\parindent}{0pt}
    \setlength{\parskip}{6pt plus 2pt minus 1pt}}
}{% if KOMA class
  \KOMAoptions{parskip=half}}
\makeatother
\usepackage{color}
\usepackage{fancyvrb}
\newcommand{\VerbBar}{|}
\newcommand{\VERB}{\Verb[commandchars=\\\{\}]}
\DefineVerbatimEnvironment{Highlighting}{Verbatim}{commandchars=\\\{\}}
% Add ',fontsize=\small' for more characters per line
\newenvironment{Shaded}{}{}
\newcommand{\AlertTok}[1]{\textcolor[rgb]{1.00,0.00,0.00}{\textbf{#1}}}
\newcommand{\AnnotationTok}[1]{\textcolor[rgb]{0.38,0.63,0.69}{\textbf{\textit{#1}}}}
\newcommand{\AttributeTok}[1]{\textcolor[rgb]{0.49,0.56,0.16}{#1}}
\newcommand{\BaseNTok}[1]{\textcolor[rgb]{0.25,0.63,0.44}{#1}}
\newcommand{\BuiltInTok}[1]{\textcolor[rgb]{0.00,0.50,0.00}{#1}}
\newcommand{\CharTok}[1]{\textcolor[rgb]{0.25,0.44,0.63}{#1}}
\newcommand{\CommentTok}[1]{\textcolor[rgb]{0.38,0.63,0.69}{\textit{#1}}}
\newcommand{\CommentVarTok}[1]{\textcolor[rgb]{0.38,0.63,0.69}{\textbf{\textit{#1}}}}
\newcommand{\ConstantTok}[1]{\textcolor[rgb]{0.53,0.00,0.00}{#1}}
\newcommand{\ControlFlowTok}[1]{\textcolor[rgb]{0.00,0.44,0.13}{\textbf{#1}}}
\newcommand{\DataTypeTok}[1]{\textcolor[rgb]{0.56,0.13,0.00}{#1}}
\newcommand{\DecValTok}[1]{\textcolor[rgb]{0.25,0.63,0.44}{#1}}
\newcommand{\DocumentationTok}[1]{\textcolor[rgb]{0.73,0.13,0.13}{\textit{#1}}}
\newcommand{\ErrorTok}[1]{\textcolor[rgb]{1.00,0.00,0.00}{\textbf{#1}}}
\newcommand{\ExtensionTok}[1]{#1}
\newcommand{\FloatTok}[1]{\textcolor[rgb]{0.25,0.63,0.44}{#1}}
\newcommand{\FunctionTok}[1]{\textcolor[rgb]{0.02,0.16,0.49}{#1}}
\newcommand{\ImportTok}[1]{\textcolor[rgb]{0.00,0.50,0.00}{\textbf{#1}}}
\newcommand{\InformationTok}[1]{\textcolor[rgb]{0.38,0.63,0.69}{\textbf{\textit{#1}}}}
\newcommand{\KeywordTok}[1]{\textcolor[rgb]{0.00,0.44,0.13}{\textbf{#1}}}
\newcommand{\NormalTok}[1]{#1}
\newcommand{\OperatorTok}[1]{\textcolor[rgb]{0.40,0.40,0.40}{#1}}
\newcommand{\OtherTok}[1]{\textcolor[rgb]{0.00,0.44,0.13}{#1}}
\newcommand{\PreprocessorTok}[1]{\textcolor[rgb]{0.74,0.48,0.00}{#1}}
\newcommand{\RegionMarkerTok}[1]{#1}
\newcommand{\SpecialCharTok}[1]{\textcolor[rgb]{0.25,0.44,0.63}{#1}}
\newcommand{\SpecialStringTok}[1]{\textcolor[rgb]{0.73,0.40,0.53}{#1}}
\newcommand{\StringTok}[1]{\textcolor[rgb]{0.25,0.44,0.63}{#1}}
\newcommand{\VariableTok}[1]{\textcolor[rgb]{0.10,0.09,0.49}{#1}}
\newcommand{\VerbatimStringTok}[1]{\textcolor[rgb]{0.25,0.44,0.63}{#1}}
\newcommand{\WarningTok}[1]{\textcolor[rgb]{0.38,0.63,0.69}{\textbf{\textit{#1}}}}
\usepackage{longtable,booktabs,array}
\newcounter{none} % for unnumbered tables
\usepackage{calc} % for calculating minipage widths
% Correct order of tables after \paragraph or \subparagraph
\usepackage{etoolbox}
\makeatletter
\patchcmd\longtable{\par}{\if@noskipsec\mbox{}\fi\par}{}{}
\makeatother
% Allow footnotes in longtable head/foot
\IfFileExists{footnotehyper.sty}{\usepackage{footnotehyper}}{\usepackage{footnote}}
\makesavenoteenv{longtable}
\ifLuaTeX
\usepackage[bidi=basic,shorthands=off]{babel}
\else
\usepackage[bidi=default,shorthands=off]{babel}
\fi
\ifLuaTeX
  \usepackage{selnolig} % disable illegal ligatures
\fi
\setlength{\emergencystretch}{3em} % prevent overfull lines
\providecommand{\tightlist}{%
  \setlength{\itemsep}{0pt}\setlength{\parskip}{0pt}}
\usepackage{bookmark}
\IfFileExists{xurl.sty}{\usepackage{xurl}}{} % add URL line breaks if available
\urlstyle{same}
\hypersetup{
  pdftitle={Compte rendu : localisation de sources par maximum de vraisemblance et critère BIC},
  pdflang={fr},
  hidelinks,
  pdfcreator={LaTeX via pandoc}}

\title{Compte rendu : localisation de sources par maximum de
vraisemblance et critère BIC}
\usepackage{etoolbox}
\makeatletter
\providecommand{\subtitle}[1]{% add subtitle to \maketitle
  \apptocmd{\@title}{\par {\large #1 \par}}{}{}
}
\makeatother
\subtitle{Méthode, résultats et limites}
\author{}
\date{}

\begin{document}
\maketitle

\section{Compte rendu : maximum de vraisemblance +
BIC}\label{compte-rendu-maximum-de-vraisemblance-bic}

Branche \texttt{travail-sur-le-probleme-inverse}, dossier
\texttt{sparse\_sources/}.

Ce compte rendu présente la méthode qui donne les meilleurs résultats
sur le problème du projet : le \textbf{maximum de vraisemblance glouton
avec choix du nombre de sources par le critère BIC} (« ML + BIC »). Il
explique l'idée, ce que la méthode permet de retrouver, et où elle
s'arrête. Les détails techniques (dérivations, tableaux complets, liste
des tests) sont dans \texttt{RAPPORT\_ML\_BIC.md} ; la comparaison avec
les autres méthodes (LASSO, BLASSO, Prony) est dans
\texttt{COMPTE\_RENDU.md}. Tous les chiffres viennent des fichiers CSV
de \texttt{sparse\_sources/results/}.

\subsection{1. Le problème}\label{le-probluxe8me}

On cherche des sources ponctuelles dans le rectangle
\(D = [0,10]\times[0,20]\) : leur nombre \(K\), leurs positions \(z_k\)
et leurs intensités \(I_k\). On dispose de \(M = 64\) capteurs placés
sur le cercle qui entoure \(D\). Le capteur \(i\), en \(x_i\), reçoit de
chaque source un signal qui décroît comme l'inverse du carré de la
distance :

\[
y_i = \sum_{k=1}^{K} \frac{I_k}{\lVert x_i - z_k\rVert^2} \;+\; \frac{I_c}{\lVert x_i - c\rVert^2} .
\]

Le second terme est la \textbf{source centrale} du projet : une source
connue, très intense (\(I_c = 5000\), contre une intensité totale de 100
pour les autres sources), placée au centre \(c\) de \(D\).

Chaque mesure est bruitée : le capteur lit
\(y_i(1 + \sigma\varepsilon_i)\), avec \(\varepsilon_i\) gaussien centré
réduit et \(\sigma\) entre \(10^{-4}\) et \(10^{-2}\). On retire ensuite
la contribution connue de la source centrale.

\textbf{Ce qui rend le problème difficile.} Retirer la source centrale
enlève son signal, mais pas son bruit. Chaque capteur reçoit
\(5000/125 = 40\) de la source centrale, contre environ 7 en moyenne
pour toutes les autres sources réunies. Le bruit restant, d'écart type
environ \(40\,\sigma\), est donc dominé par la source centrale.

\subsection{2. L'idée : deux questions
séparées}\label{liduxe9e-deux-questions-suxe9paruxe9es}

La méthode répond séparément à deux questions.

\textbf{Où sont les sources et quelle est leur intensité ? Le maximum de
vraisemblance (ML).} Pour un nombre de sources \(K\) donné, on cherche
les positions et intensités qui rendent les mesures observées
\textbf{les plus probables}. Avec un bruit gaussien, cela revient à
minimiser l'écart entre les mesures \(m_i\) et le modèle, chaque capteur
étant divisé par l'écart type \(s_i\) de son bruit :

\[
\chi^2 = \sum_{i=1}^{M}\left(\frac{m_i - \text{modèle}_i}{s_i}\right)^2 .
\]

Diviser par \(s_i\) revient à faire moins confiance aux capteurs les
plus bruités. On ajuste positions et intensités \textbf{sans aucune
pénalité} : il n'y a pas de paramètre de régularisation qui fausserait
le résultat.

\textbf{Combien y a-t-il de sources ? Le critère BIC.} Le ML seul ne
suffit pas : plus on ajoute de sources, plus \(\chi^2\) baisse, même si
les sources ajoutées sont fausses, car une source de trop « explique »
toujours un peu de bruit. Le critère d'information bayésien (BIC) fait
payer chaque source :

\[
C(K) = \chi^2 + 3K\log M .
\]

Chaque source a 3 paramètres (\(x\), \(y\), \(I\)) et chaque paramètre
coûte \(\log M\approx 4{,}2\) pour 64 capteurs. \textbf{Une source n'est
donc gardée que si elle fait baisser \(\chi^2\) de plus de
\(3\log 64\approx 12{,}5\).}

Pourquoi ce seuil fonctionne :

\begin{itemize}
\tightlist
\item
  une \textbf{fausse source}, qui n'explique que du bruit, fait baisser
  \(\chi^2\) d'environ 3 en moyenne (un par paramètre) : elle est
  rejetée ;
\item
  une \textbf{vraie source} fait baisser \(\chi^2\) d'environ
  \(\mathrm{SNR}^2 + 3\), où \(\mathrm{SNR}\) est son rapport signal sur
  bruit : elle est gardée dès que \(\mathrm{SNR}\gtrsim 3\).
\end{itemize}

Il n'y a \textbf{aucun réglage} : le seuil ne dépend que du nombre de
capteurs.

\textbf{Comparaison avec le BLASSO.} Le BLASSO (Frank-Wolfe) règle ces
deux questions avec un seul paramètre \(\lambda\) : un \(\lambda\) petit
laisse passer de fausses sources, un \(\lambda\) grand décale les
positions et sous-estime les intensités. Il ne fonctionne en plus que si
une condition mathématique (la condition de Duval--Peyré) est vraie, ce
qui est rare à partir de 3 sources. ML + BIC n'a aucune de ces deux
contraintes.

\subsection{3. L'algorithme}\label{lalgorithme}

On construit la solution source par source (algorithme « glouton ») :

\begin{enumerate}
\def\labelenumi{\arabic{enumi}.}
\tightlist
\item
  on part de zéro source ;
\item
  on cherche l'endroit où les mesures encore inexpliquées ressemblent le
  plus au signal d'une source, et on y place une nouvelle source ;
\item
  on réajuste \textbf{ensemble} toutes les positions et intensités pour
  minimiser \(\chi^2\) (optimisation L-BFGS-B, gradient exact) ;
\item
  on calcule \(C(K)\) ; s'il baisse, on recommence en 2 ; s'il remonte
  deux fois de suite, on s'arrête ;
\item
  étape arrière : on essaie de retirer chaque source, et on la retire si
  \(C\) baisse ;
\item
  on renvoie la solution de plus petit \(C(K)\).
\end{enumerate}

Exemple réel (3 sources, \(\sigma = 10^{-3}\), avec la source centrale)
:

{\def\LTcaptype{none} % do not increment counter
\begin{longtable}[]{@{}lllllll@{}}
\toprule\noalign{}
\(K\) & 0 & 1 & 2 & \textbf{3} & 4 & 5 \\
\midrule\noalign{}
\endhead
\bottomrule\noalign{}
\endlastfoot
\(C(K)\) & 735 305 & 238 756 & 997 & \textbf{89,8} & 97,7 & 102,9 \\
\end{longtable}
}

Le critère chute jusqu'à \(K = 3\), puis remonte : la 4e source fait
baisser \(\chi^2\) de moins de 12,5. La méthode trouve 3 sources. Elle
prend 0,1 à 0,5 seconde pour 1 à 5 sources.

\subsection{4. Savoir si le résultat est
fiable}\label{savoir-si-le-ruxe9sultat-est-fiable}

La méthode fournit deux contrôles.

\textbf{Le \(\chi^2\) final.} Si le modèle explique les mesures,
\(\chi^2/M\) vaut environ \((M - 3K)/M\), soit 0,86 pour 3 sources et 64
capteurs. Une valeur nettement supérieure à 1 signale une source manquée
ou un bruit sous-estimé.

\textbf{La borne de Cramér--Rao (CRB).} C'est l'erreur de position
minimale que \textbf{n'importe quelle méthode} peut atteindre en
moyenne, compte tenu du bruit. Elle se calcule pour chaque source
(\texttt{cramer\_rao}). Si la CRB d'une source est grande devant la
distance aux autres sources, cette source ne peut pas être localisée de
façon fiable, quel que soit l'algorithme. C'est l'outil qui permet de
distinguer « la méthode est mauvaise » de « le problème est impossible
».

\subsection{5. Résultats sur le scénario du
projet}\label{ruxe9sultats-sur-le-scuxe9nario-du-projet}

Mêmes données que \texttt{scripts/reconstruct.py}, 20 tirages par case.
Une reconstruction est \textbf{réussie} si le nombre de sources est
exact et si chaque source est retrouvée à moins d'un quart de la
distance entre les deux sources les plus proches.

{\def\LTcaptype{none} % do not increment counter
\begin{longtable}[]{@{}lllll@{}}
\toprule\noalign{}
sources & bruit \(\sigma\) & Prony & Frank-Wolfe & \textbf{ML + BIC} \\
\midrule\noalign{}
\endhead
\bottomrule\noalign{}
\endlastfoot
1 & \(10^{-3}\) & 15 \% & 100 \% & \textbf{100 \%} \\
1 & \(10^{-2}\) & 0 \% & 100 \% & \textbf{100 \%} \\
2 & \(10^{-4}\) & 85 \% & 90 \% & \textbf{95 \%} \\
2 & \(10^{-3}\) & 15 \% & 70 \% & \textbf{75 \%} \\
2 & \(3\times 10^{-3}\) & 0 \% & 40 \% & \textbf{75 \%} \\
3 & \(10^{-4}\) & 55 \% & 40 \% & \textbf{80 \%} \\
3 & \(10^{-3}\) & 10 \% & 5 \% & \textbf{40 \%} \\
5 & \(10^{-4}\) & 5 \% & 0 \% & \textbf{50 \%} \\
\end{longtable}
}

\begin{enumerate}
\def\labelenumi{\arabic{enumi}.}
\tightlist
\item
  \textbf{ML + BIC est la meilleure méthode dans toutes les cases}, et
  l'écart grandit avec le nombre de sources et le bruit.
\item
  \textbf{Elle est optimale} : son erreur vaut 0,7 à 0,9 fois la CRB. On
  ne peut donc pas faire beaucoup mieux en changeant d'algorithme.
\item
  \textbf{Ses échecs sont des cas impossibles} : sur les tirages où la
  CRB permet de réussir, elle réussit dans 90 à 100 \% des cas.
\end{enumerate}

\textbf{Le vrai frein est la source centrale.} Pour 3 sources à
\(\sigma = 10^{-3}\), ML + BIC réussit dans \textbf{100 \%} des cas sans
source centrale et dans \textbf{40 \%} avec. La source centrale
multiplie la CRB par environ 42.

\subsection{6. Ce que la méthode peut
voir}\label{ce-que-la-muxe9thode-peut-voir}

\textbf{Une source faible} (une source seule, \(\sigma = 10^{-3}\),
source centrale) :

{\def\LTcaptype{none} % do not increment counter
\begin{longtable}[]{@{}lllllll@{}}
\toprule\noalign{}
intensité & 0,5 & 1 & 2 & 5 & 10 & 20 \\
\midrule\noalign{}
\endhead
\bottomrule\noalign{}
\endlastfoot
détectée & 25 \% & 48 \% & 90 \% & 100 \% & 100 \% & 100 \% \\
localisée à moins de 1 & 18 \% & 22 \% & 40 \% & 80 \% & 98 \% & 100
\% \\
\end{longtable}
}

La probabilité de détection suit exactement la prédiction théorique du
BIC (seuil \(\mathrm{SNR}\approx 3\)). Sur du bruit seul, la méthode
invente une source dans 2 \% des cas seulement. Une source d'intensité 2
est presque toujours vue, mais il en faut environ 10 pour la placer à
mieux que 1 près.

\textbf{Deux sources proches} (intensité 50 chacune) : elles sont
séparées à partir d'une distance d'environ 1 à \(\sigma = 10^{-5}\), 2 à
\(\sigma = 10^{-4}\) et 4 à \(\sigma = 10^{-3}\). En dessous, la méthode
trouve une seule source au lieu de deux, ce qui vaut mieux que d'en
inventer.

\subsection{7. Robustesse}\label{robustesse}

\textbf{Bruit mal connu.} La méthode a besoin du niveau de bruit
\(\sigma\). Le \textbf{surestimer} d'un facteur 2 ne coûte presque rien
(on perd seulement les sources les plus faibles). Le
\textbf{sous-estimer} d'un facteur 2 fait apparaître en moyenne une
source parasite. En cas de doute, il faut donc prendre \(\sigma\) un peu
trop grand.

\textbf{Plus de capteurs sur le même cercle.} La CRB baisse comme
\(1/\sqrt M\) : quadrupler le nombre de capteurs divise l'erreur par 2.
Pour 3 sources à \(\sigma = 10^{-3}\), la réussite passe de 40 \% (64
capteurs) à 70 \% (1024 capteurs). Le gain est réel mais lent.

\textbf{Capteurs en hauteur, sur la paroi d'un cylindre.} Pas de gain :
avec 3 sources, le cercle simple reste le meilleur, et à partir de 5
sources toutes les géométries échouent, pour le Frank-Wolfe comme pour
le ML.

\subsection{8. En 3D}\label{en-3d}

La méthode se généralise sans modification à des sources dans un cube
(\texttt{ml3d.py}) : chaque source a alors 4 paramètres, d'où un seuil
de \(4\log M\). Sur le cube \([0,10]^3\) avec 160 capteurs et la source
centrale :

\begin{itemize}
\tightlist
\item
  avec des capteurs sur une \textbf{sphère}, 2 sources sont retrouvées
  dans 100 \% des cas à \(\sigma = 10^{-4}\) et 88 \% à
  \(\sigma = 10^{-3}\), 5 sources dans 50 \% des cas à
  \(\sigma = 10^{-4}\) ;
\item
  l'erreur vaut 0,75 à 0,9 fois la CRB, comme en 2D ;
\item
  le BLASSO 3D du projet échoue dans presque tous les cas ;
\item
  \textbf{un seul cercle de capteurs ne suffit pas en 3D} : une source
  au-dessus du cercle et sa symétrique en dessous donnent les mêmes
  mesures. Il faut des capteurs hors d'un même plan (plusieurs cercles,
  ou une sphère, qui est la meilleure disposition).
\end{itemize}

\subsection{9. Et avec beaucoup de sources
?}\label{et-avec-beaucoup-de-sources}

Expérience \texttt{beaucoup} de \texttt{experiments\_ml.py}
(\texttt{ml\_beaucoup\_directions.csv}, \texttt{ml\_beaucoup\_ml.csv},
\texttt{ml\_beaucoup.png}, 5 tirages par cas). Sources tirées comme dans
\texttt{config/} : positions uniformes, intensités uniformes sur le
simplexe de somme 100.

\textbf{Ce que les mesures contiennent.} Avec \(K\) sources, il y a
\(3K\) paramètres à retrouver. On compte combien de « directions » de
cet espace de paramètres les capteurs mesurent vraiment : ce sont les
valeurs singulières de la matrice des dérivées du modèle (colonnes
normalisées) qui dépassent \(10^{-6}\) fois la plus grande. Le calcul
est fait \textbf{sans bruit} : un bruit réaliste (\(\sigma = 10^{-3}\))
en laisserait encore moins.

{\def\LTcaptype{none} % do not increment counter
\begin{longtable}[]{@{}
  >{\raggedright\arraybackslash}p{(\linewidth - 8\tabcolsep) * \real{0.2000}}
  >{\raggedright\arraybackslash}p{(\linewidth - 8\tabcolsep) * \real{0.2000}}
  >{\raggedright\arraybackslash}p{(\linewidth - 8\tabcolsep) * \real{0.2000}}
  >{\raggedright\arraybackslash}p{(\linewidth - 8\tabcolsep) * \real{0.2000}}
  >{\raggedright\arraybackslash}p{(\linewidth - 8\tabcolsep) * \real{0.2000}}@{}}
\toprule\noalign{}
\begin{minipage}[b]{\linewidth}\raggedright
sources \(K\)
\end{minipage} & \begin{minipage}[b]{\linewidth}\raggedright
paramètres \(3K\)
\end{minipage} & \begin{minipage}[b]{\linewidth}\raggedright
directions mesurables, 64 capteurs
\end{minipage} & \begin{minipage}[b]{\linewidth}\raggedright
1024 capteurs
\end{minipage} & \begin{minipage}[b]{\linewidth}\raggedright
plus petite distance entre sources
\end{minipage} \\
\midrule\noalign{}
\endhead
\bottomrule\noalign{}
\endlastfoot
1 & 3 & 3 & 3 & \\
3 & 9 & 9 & 9 & 4,3 \\
5 & 15 & 15 & 15 & 2,7 \\
10 & 30 & 29 & 29 & 0,74 \\
20 & 60 & 46 & 47 & 0,62 \\
30 & 90 & 50 & 56 & 0,24 \\
50 & 150 & 56 & 71 & 0,15 \\
\end{longtable}
}

\begin{itemize}
\tightlist
\item
  Jusqu'à 5 sources, tous les paramètres sont mesurés. Au-delà, le
  nombre de directions mesurables \textbf{plafonne autour de 50 à 70},
  même avec 16 fois plus de capteurs.
\item
  La raison : vu depuis le cercle, le signal en \(1/r^2\) est très
  lisse. Les capteurs voient bien les grandes structures (intensité
  totale, barycentre, étalement), mais les détails fins sont écrasés.
  Des capteurs supplémentaires sur le même cercle répètent presque la
  même information.
\item
  Avec 50 sources, il y a de plus 150 inconnues pour 64 mesures : le
  problème n'a pas de solution unique, même sans bruit. Et les sources
  tirées au hasard sont à environ 0,15 les unes des autres, alors que
  même deux sources isolées ne sont séparées qu'à partir d'une distance
  de 2 à 4 avec le bruit du projet.
\end{itemize}

\textbf{Ce que ML + BIC retrouve} (64 capteurs, source centrale, médiane
sur 5 tirages) :

{\def\LTcaptype{none} % do not increment counter
\begin{longtable}[]{@{}
  >{\raggedright\arraybackslash}p{(\linewidth - 8\tabcolsep) * \real{0.2000}}
  >{\raggedright\arraybackslash}p{(\linewidth - 8\tabcolsep) * \real{0.2000}}
  >{\raggedright\arraybackslash}p{(\linewidth - 8\tabcolsep) * \real{0.2000}}
  >{\raggedright\arraybackslash}p{(\linewidth - 8\tabcolsep) * \real{0.2000}}
  >{\raggedright\arraybackslash}p{(\linewidth - 8\tabcolsep) * \real{0.2000}}@{}}
\toprule\noalign{}
\begin{minipage}[b]{\linewidth}\raggedright
vraies sources
\end{minipage} & \begin{minipage}[b]{\linewidth}\raggedright
\(\sigma\)
\end{minipage} & \begin{minipage}[b]{\linewidth}\raggedright
sources estimées
\end{minipage} & \begin{minipage}[b]{\linewidth}\raggedright
vraies sources retrouvées à moins de 0,5
\end{minipage} & \begin{minipage}[b]{\linewidth}\raggedright
temps
\end{minipage} \\
\midrule\noalign{}
\endhead
\bottomrule\noalign{}
\endlastfoot
10 & \(10^{-4}\) & 6 & 5 & 1,2 s \\
10 & \(10^{-3}\) & 4 & 2 & 0,3 s \\
20 & \(10^{-4}\) & 8 & 4 & 4,6 s \\
20 & \(10^{-3}\) & 6 & 3 & 0,9 s \\
50 & \(10^{-4}\) & 10 & 6 & 14 s \\
50 & \(10^{-3}\) & 7 & 4 & 1,5 s \\
\end{longtable}
}

La méthode ne s'effondre pas : elle garde \textbf{les quelques sources
que les données permettent de voir} (les plus intenses et les plus
isolées) et n'invente pas les autres. Mais \textbf{aucune méthode ne
peut retrouver 50 sources} avec ces mesures. Pour s'en approcher, il
faudrait changer les mesures : capteurs proches des sources (à
l'intérieur de \(D\)), bruit beaucoup plus faible, ou information a
priori (par exemple des positions possibles connues, ce qui ne laisse
que les intensités à estimer). On peut aussi viser des grandeurs
globales (intensité totale, barycentre, carte de densité grossière)
plutôt que chaque source.

\subsection{10. Conclusion}\label{conclusion}

{\def\LTcaptype{none} % do not increment counter
\begin{longtable}[]{@{}
  >{\raggedright\arraybackslash}p{(\linewidth - 2\tabcolsep) * \real{0.5000}}
  >{\raggedright\arraybackslash}p{(\linewidth - 2\tabcolsep) * \real{0.5000}}@{}}
\toprule\noalign{}
\begin{minipage}[b]{\linewidth}\raggedright
situation (2D, source centrale du projet)
\end{minipage} & \begin{minipage}[b]{\linewidth}\raggedright
ce que ML + BIC retrouve
\end{minipage} \\
\midrule\noalign{}
\endhead
\bottomrule\noalign{}
\endlastfoot
\(\sigma = 10^{-4}\) & 1 à 3 sources de façon fiable, 5 sources une fois
sur deux \\
\(\sigma = 10^{-3}\) & 1 à 2 sources, 3 sources dans 40 \% des cas \\
\(\sigma = 10^{-2}\) & 1 source \\
20 à 50 sources & seulement les quelques sources les plus intenses
(section 9) \\
\end{longtable}
}

\begin{enumerate}
\def\labelenumi{\arabic{enumi}.}
\tightlist
\item
  \textbf{ML + BIC est la méthode à utiliser} : la meilleure dans tous
  les tests, sans paramètre à régler, et au niveau de la borne de
  Cramér--Rao.
\item
  \textbf{Ses limites sont celles des données}, pas de l'algorithme :
  bruit de la source centrale, nombre de sources, sources trop faibles
  ou trop proches.
\item
  \textbf{Pour aller plus loin, il faut changer les mesures} : réduire
  le bruit de la source centrale (c'est le facteur dominant), rapprocher
  des capteurs des sources, ou ajouter de l'information a priori
  (positions possibles connues, par exemple). La CRB permet d'évaluer
  chaque option avant de la mettre en œuvre.
\end{enumerate}

\subsection{11. Utilisation}\label{utilisation}

\begin{Shaded}
\begin{Highlighting}[]
\ImportTok{from}\NormalTok{ sparse\_sources }\ImportTok{import}\NormalTok{ greedy\_ml, cramer\_rao, relative\_noise\_std, make\_setup}

\NormalTok{S }\OperatorTok{=}\NormalTok{ make\_setup(h}\OperatorTok{=}\FloatTok{0.0}\NormalTok{, normalize}\OperatorTok{=}\VariableTok{True}\NormalTok{)       }\CommentTok{\# capteurs du projet}
\NormalTok{s }\OperatorTok{=}\NormalTok{ relative\_noise\_std(y\_brut, sigma)       }\CommentTok{\# bruit de chaque capteur (source centrale comprise)}
\NormalTok{r }\OperatorTok{=}\NormalTok{ greedy\_ml(m, s, setup}\OperatorTok{=}\NormalTok{S)                }\CommentTok{\# m : mesures, source centrale retirée}
\NormalTok{r[}\StringTok{"z"}\NormalTok{], r[}\StringTok{"I"}\NormalTok{], r[}\StringTok{"K"}\NormalTok{]                      }\CommentTok{\# positions, intensités, nombre de sources}
\NormalTok{r[}\StringTok{"chi2"}\NormalTok{] }\OperatorTok{/}\NormalTok{ S.M                             }\CommentTok{\# contrôle : environ (M {-} 3K) / M}
\NormalTok{cramer\_rao(r[}\StringTok{"z"}\NormalTok{], r[}\StringTok{"I"}\NormalTok{], s, S)[}\StringTok{"pos\_std"}\NormalTok{] }\CommentTok{\# précision de chaque source}
\end{Highlighting}
\end{Shaded}

Pour tester la méthode sur un tirage, avec une figure :

\begin{Shaded}
\begin{Highlighting}[]
\ExtensionTok{python}\NormalTok{ sparse\_sources/tester\_ml.py                       }\CommentTok{\# 3 sources, sigma = 1e{-}3}
\ExtensionTok{python}\NormalTok{ sparse\_sources/tester\_ml.py }\AttributeTok{{-}{-}sources}\NormalTok{ 5 }\AttributeTok{{-}{-}sigma}\NormalTok{ 1e{-}4 }\AttributeTok{{-}{-}comparer}
\ExtensionTok{python}\NormalTok{ sparse\_sources/tester\_ml.py }\AttributeTok{{-}{-}dim}\NormalTok{ 3 }\AttributeTok{{-}{-}capteurs}\NormalTok{ sphere}
\end{Highlighting}
\end{Shaded}

Pour reproduire les résultats :

\begin{Shaded}
\begin{Highlighting}[]
\ExtensionTok{pytest}\NormalTok{ sparse\_sources/tests                    }\CommentTok{\# 86 tests}
\ExtensionTok{python}\NormalTok{ sparse\_sources/experiments\_projet.py    }\CommentTok{\# comparaison des méthodes (section 5)}
\ExtensionTok{python}\NormalTok{ sparse\_sources/experiments\_ml.py        }\CommentTok{\# sections 6, 7 et 9}
\ExtensionTok{python}\NormalTok{ sparse\_sources/experiments\_cylindre.py  }\CommentTok{\# capteurs sur un cylindre}
\ExtensionTok{python}\NormalTok{ sparse\_sources/experiments\_ml3d.py      }\CommentTok{\# 3D (section 8)}
\end{Highlighting}
\end{Shaded}

\subsection{Références}\label{ruxe9fuxe9rences}

\begin{itemize}
\tightlist
\item
  G. Schwarz, \emph{Estimating the dimension of a model}, The Annals of
  Statistics, 1978 (critère BIC).
\item
  S. M. Kay, \emph{Fundamentals of Statistical Signal Processing},
  volumes I (estimation, borne de Cramér--Rao) et II (détection),
  Prentice Hall, 1993 et 1998.
\item
  V. Duval, G. Peyré, \emph{Exact support recovery for sparse spikes
  deconvolution}, Foundations of Computational Mathematics, 2015
  (BLASSO).
\end{itemize}

\end{document}
